\begin{frame}{할당 단계가 $J$를 줄이는 이유}
  \begin{itemize}
    \item $J$는 각 점의 기여 $\|x - \mu_{\text{own}}\|^2$의 합
    \item 어떤 점이 자기 중심점에서 \textbf{다른} 중심점보다
      가까워졌다면, 그 점만 그 그룹으로 옮기면
      \textbf{해당 기여가 작아짐}
    \item $\Rightarrow$ 전체 합 $J$는 \textbf{줄거나 그대로}
  \end{itemize}
  \vskip0.8em
  \begin{block}{핵심}
    ``모든 점을 \textbf{가장 가까운} 중심점에게 배정한다''는
    규칙은, 정확히 \textbf{이 기여를 각 점마다 최소화}하는
    동작이다.
  \end{block}
\end{frame}
